\section{Product tests and regular trajectories}
\label{sec:framework}

We work on the Hilbert space $(\C^2)^{\otimes m}$.  Vector norms are
Euclidean norms, and operator norms are the induced norms.  Unless explicitly
stated otherwise, every starting vector has norm at most one.  Products of
orthogonal projections used in the analysis are applied without renormalizing
the resulting vectors.
All projection insertions below are analytic devices; the permitted
circuit operations remain unitary until the final output measurement.

\subsection{The gate model and product tests}

The physical circuit has $n$ input qubits, $a$ ancillary qubits initialized to
$\lvert0\rangle^{\otimes a}$, and a designated output qubit measured in the
computational basis.  Its gates are arbitrary one-qubit unitaries and the
generalized Toffoli gates of Section~\ref{sec:introduction}.  Gates in each layer have pairwise
disjoint supports, and each gate counts once regardless of its arity.
The gates may depend arbitrarily on the input length,
including through their complex entries, but are independent of the input
string.  No input or ancillary register is required to be restored, and the
output qubit may be an original input qubit.  All other qubits may be discarded
after the output is measured.  There are no intermediate measurements,
postselection operations, geometric restrictions, or extra fanout gates.

For the analysis, we allow a reflection
\[
 R_g=I-2P_g,
 \qquad
 P_g=\bigotimes_{j\in J_g}\lvert v_{g,j}\rangle\langle v_{g,j}\rvert
       \otimes I_{J_g^c},
\]
where $J_g$ is a nonempty set of sites and each $v_{g,j}$ is a unit vector.
The tensor factors are understood to occupy their indicated sites.  Each
$P_g$ is an orthogonal projection, so $R_g=R_g^\dagger=R_g^{-1}$.
A generalized Toffoli with $c\ge1$ distinct controls $z_1,\ldots,z_c$ and
a distinct target $b$ has this form with
\[
 P_g=\left(\bigotimes_{i=1}^{c}\lvert1\rangle\langle1\rvert_{z_i}\right)
       \otimes\lvert-\rangle\langle-\rvert_b,
 \qquad
 \lvert-\rangle=\frac{\lvert0\rangle-\lvert1\rangle}{\sqrt2},
\]
and identities elsewhere.  Indeed, it acts trivially unless all controls are
one, and on that control subspace its target action is
$I-2\lvert-\rangle\langle-\rvert=X$.
This product-reflection viewpoint also appears in
\citet{Rosenthal2021}.

An \emph{analytic layer} will be either a product of such reflections on
pairwise disjoint supports, or a tensor product of one-qubit unitaries.
Splitting a physical layer into its two gate types preserves its operation,
because its gates have disjoint supports.  Thus a physical circuit of depth
$d$ gives at most $2d$ analytic layers.  We use $t$ for their number and $m$
for the total number of qubits, including any reference qubits subsequently
introduced.  Empty layers are permitted.  Every reflection layer has at most
$m$ gates, since their supports are nonempty and disjoint.

\begin{definition}[Product test]
\label{def:test}
A test $A$ consists of a subset $J_A\subseteq\{1,\ldots,m\}$ and a specified
unit vector $v_{A,j}$ at each $j\in J_A$.  Its excitation count and zero
projection are
\[
 H_A=\sum_{j\in J_A}
       \bigl(I-\lvert v_{A,j}\rangle\langle v_{A,j}\rvert\bigr)_j,
 \qquad
 P_A=\one_{H_A=0}
    =\bigotimes_{j\in J_A}\lvert v_{A,j}\rangle\langle v_{A,j}\rvert
       \otimes I_{J_A^c}.
\]
Here a subscript $j$ extends a one-site operator by identity at all other
sites, and $\one_{H_A\in E}$ denotes the spectral projection of $H_A$ onto
the set $E\subseteq\mathbb R$.  The empty test has $H_A=0$ and $P_A=I$.
\end{definition}

The summands of $H_A$ are commuting orthogonal projections.  Choosing the
basis $(v_{A,j},v_{A,j}^{\perp})$ at each tested site diagonalizes them
simultaneously.  Consequently its spectrum is contained in
$\{0,1,\ldots,|J_A|\}$, and $0\le H_A\le mI$.  We use spectral inequalities
such as $H_A\ge x$ also when $x$ is not an integer.

For a set $S$ of gates in one reflection layer, disjointness makes their
specified local factors a single test.  We write $H_S$ for its count and
$P_S=\prod_{g\in S}P_g$ for its zero projection.  If $F$ is the whole
reflection layer, then
\[
 [F,H_S]=0.
\]
Indeed, a gate in $S$ commutes with each of its own local factors, while
other gates have disjoint supports.  Every spectral projection of $H_S$
therefore commutes with $F$ as well.  If $H_A$ retains only some sites of
this test, then $H_A$ and $H_S$ are diagonal in the same product basis,
$H_A\le H_S$, and
\[
 \one_{H_A\ge x}\le\one_{H_S\ge x}
 \quad\text{for every real }x.
\]

\subsection{Asymptotic sequences and uniformity}

All asymptotic statements concern an arbitrary sequence indexed by
$\nu\in\mathbb N$, with total qubit counts $m_\nu\to\infty$ and a bounded
number of analytic layers.  The states, tests, gates, and circuits can all
vary with $\nu$.  We suppress $\nu$ and write $m$ for $m_\nu$; thus
``eventually'' means at all sufficiently large sequence indices.  Every
exponent appearing in an estimate is a fixed real number independent of
$\nu$.  Padding with empty layers gives a common bound on length.  When a
fixed layer pattern is convenient, one may restrict to a subsequence with
fixed length and layer types: there are only finitely many such patterns.
An eventual assertion on every such subsequence yields the assertion on
the original sequence, since any infinite set of violations has an infinite
subsequence with one fixed pattern.

The following elementary observation will also supply uniformity over
tests chosen from the current state.

\begin{lemma}[Uniformity from arbitrary choices]
\label{lem:uniformization}
Fix a sequence of base data, and let $\mathcal D_\nu$ be the nonempty set of
permitted auxiliary choices at index $\nu$.  Suppose that for every sequence
$D_\nu\in\mathcal D_\nu$ and every fixed admissible parameter tuple
$\theta$, an assertion $\mathcal E(\nu,D_\nu,\theta)$ holds eventually.
Then, for every fixed $\theta$, there is an index $\nu_0$ such that
\[
 \mathcal E(\nu,D,\theta)
 \quad\text{holds for all }\nu\ge\nu_0
 \text{ and every }D\in\mathcal D_\nu.
\]
The same conclusion holds if the premise is proved on every subsequence
on which the base hypotheses remain valid.
\end{lemma}

\begin{proof}
Fix $\theta$.  If the conclusion fails, there are increasing indices
$\nu_i$ and choices $D_{\nu_i}\in\mathcal D_{\nu_i}$ violating the
assertion.  Choose arbitrary permitted values at the other indices.  The
resulting sequence of choices contradicts the premise.  In the subsequence
version, apply the premise directly to the violating subsequence.
\end{proof}

The lemma imposes no cardinality or finiteness requirement on
$\mathcal D_\nu$.  In particular, a permitted test may depend on the exact
state at its index.  The eventual threshold may depend on the fixed tuple
$\theta$; no threshold uniform over all real exponent tuples is asserted.

\begin{definition}[Regular trajectory]
\label{def:regular}
Write $V=U_t\cdots U_1$ for analytic layers and set
$\xi_j=U_j\cdots U_1\xi_0$, where $\norm{\xi_0}\le1$.
The sequence of these trajectories is \emph{regular} if the following
holds at every reflection layer $U_j$: for \emph{every fixed} triple
\[
 0\le b<h<1,\qquad 0<p<h,
\]
eventually, simultaneously for every set $S$ of at most $m^b$ gates of
that layer,
\begin{equation}
 \norm{\one_{H_S\ge m^h}\xi_{j-1}}
 \le \exp(-m^p).
 \label{eq:regularity}
\end{equation}
There is no regularity condition at one-qubit layers.  Regularity is a
property of an asymptotic sequence, not of a single finite circuit.
\end{definition}

Commutation with $H_S$ gives
\[
 \norm{\one_{H_S\ge m^h}\xi_j}
 =\norm{U_j\one_{H_S\ge m^h}\xi_{j-1}}
 =\norm{\one_{H_S\ge m^h}\xi_{j-1}}
\]
at a reflection layer, so either endpoint can be used in
\eqref{eq:regularity}.  It follows that reversing a regular trajectory is
regular for the inverse circuit: inverse reflection layers are the same
layers, and inverse one-qubit layers impose no condition.  Restricting a
regular trajectory to any segment or subsequence preserves regularity.
Multiplying each trajectory by a scalar of modulus at most one also
preserves it.  For any fixed use of regularity, finitely many auxiliary
exponents and layers may be handled simultaneously by taking the maximum
of their eventual thresholds.  This does not require a common threshold
for all admissible exponents.

\subsection{Operator degree and excitation shifts}

\begin{definition}[Operator degree]
\label{def:degree}
An operator is supported on $J\subseteq\{1,\ldots,m\}$ if it has the form
$E_J\otimes I_{J^c}$.  It has degree at most the nonnegative integer $\ell$
if it is a finite sum of operators each supported on at most $\ell$ sites.
An assertion $\deg E\le x$ for real $x\ge0$ means
$\deg E\le\lfloor x\rfloor$.
\end{definition}

Degree bounds survive sums, while products have degree at most the sum
of the factors' degrees: expand into supported terms and take the union
of each pair of supports.  Taking adjoints preserves degree, as does
conjugating by a tensor product of one-qubit unitaries.  No bound on the
number or norms of terms in a support decomposition is part of the
definition.

The following support constraint is the elementary locality principle
behind the connection between polynomial approximation and Hamming
concentration; compare \citet[Lemmas~2.10--2.11]{AnshuMetger2023}.
We record its exact form for arbitrary product tests.

\begin{lemma}[Degree controls excitation shifts]
\label{lem:degree-shift}
Let $\ell$ be a nonnegative integer.  If $\deg E\le\ell$ and $A$ is any
test, then
\[
 \one_{H_A=f}\,E\,\one_{H_A=e}=0
 \qquad\text{whenever }|f-e|>\ell.
\]
In particular, for every real $x$,
\[
 \one_{H_A>x+\ell}\,E\,\one_{H_A\le x}=0,
 \qquad
 \one_{H_A>\ell}\,EP_A=0.
\]
If $Q$ is a polynomial of degree at most $\ell$, then
$\deg Q(H_A)\le\ell$.
\end{lemma}

\begin{proof}
First let $E$ be supported on a set $J$ of size at most $\ell$.
Decompose $H_A$ into the counts on $J_A\cap J$ and on $J_A\setminus J$.
In the product basis for the test, a matrix element of $E$ can be nonzero
only when the states outside $J$ agree, so their excitation count is
unchanged.  Each count inside $J$ is between zero and $|J_A\cap J|$.
The total counts can therefore differ by at most $|J|\le\ell$.  This
proves the block assertion for a supported term, hence for their sum.
Summing the forbidden blocks gives the two displayed consequences.

Finally, expanding $H_A^j$ produces sums of products of $j$ one-site
operators.  Each product is supported on at most $j$ sites, even when a
site occurs more than once.  Expanding $Q$ into powers of $H_A$ proves the
last assertion, including its constant term.
\end{proof}
